\begin{viuabstract}{Resumen}
El transporte cooperativo de cargas heterogéneas exige formar una coalición adecuada,
ejecutar la interacción física y coordinar el tráfico con otras misiones. Este TFM
organiza el problema en tres subproblemas. SP1 decide qué robots atienden cada carga,
cuántos se necesitan, qué rol o contacto asume cada uno y cuándo esperar, iniciar o
cambiar de coalición. SP2 aborda la aproximación, el acoplamiento, la estabilización,
el transporte y la sustitución de miembros. SP3 coordina rutas, reservas y prioridades
entre robots libres, equipos en reunión y coaliciones en movimiento. \emph{Cargo}
modela una carga soportada como cuerpo rígido planar.

La propuesta combina referencias LP/MILP, dinámica primal--dual central y
vecinal, cierre entero con reparación y certificados separados para recursos,
contacto y movimiento. La evaluación usa mundos pareados, semillas explícitas y
oráculos. La batería separa siete preguntas: caso manual (E0),
relajación (E1), conectividad (E2), discretización (E3), recuperación entera (E4),
escala (E5) y llegada uniciclo (E6).

E0 reprodujo el LP regularizado en un caso controlado con errores del orden de
\(10^{-6}\)--\(10^{-5}\). En E1, Rep-D alcanzó el criterio relativo,
pero solo una ejecución fue comparable al LP bajo tolerancia primal estricta.
E2 y E3 siguen no concluyentes. E4 alcanzó factibilidad en los ocho mundos piloto y una brecha mediana
del \(1{,}21\%\). En E5, Rep-D+repair recuperó factibilidad en el caso \(N=40\),
aunque la dinámica fraccionaria no convergió dentro del presupuesto; el baseline
Greedy dejó un caso no factible. E6 obtuvo llegada completa en dos escenarios, sin
validar transporte físico.

El estado de evidencia es C-pilot: los resultados verifican implementación y casos
controlados, pero no demuestran convergencia global, optimalidad ni seguridad
distribuida del sistema completo. Las campañas de SP2 y SP3 aportan evidencia
parcial; permanecen pendientes la validación confirmatoria, la dinámica tridimensional,
la tracción rueda--suelo, las redes conmutadas, el hardware y el transporte por
\emph{caging}.

\keywords{Palabras clave}{coordinación multi-robot,
  robots móviles autónomos, formación de coaliciones,
  juegos potenciales, transporte cooperativo}
\end{viuabstract}

\clearpage
